\chapter{Cyclotomic Fields and Fermat}
\begin{center}\itshape
``It is impossible to separate a cube into two cubes, or a fourth power into two fourth powers, or in general any power
higher than the second into two like powers. I have discovered a truly marvellous proof of this, which this margin is
too narrow to contain.''\\[0.3em]
\upshape Pierre de Fermat, in the margin of his copy of Diophantus, around 1637
\end{center}

Fermat never wrote the proof down, and probably never had one. What he left was a dare, and for two centuries the best
mathematicians took it one exponent at a time. Fermat himself handled $n=4$ by infinite descent. Euler did $n=3$ in 1770,
with a gap that only the arithmetic of $\Z[\omega]$ could close (Tome~1, Theorem~10.8). Then came a surprise. Twenty years
earlier Gauss had been corresponding with a brilliant young ``Monsieur Le Blanc'', and was astonished to learn in 1807 that
Le Blanc was Sophie Germain, who had borrowed a man's name to be taken seriously. Around 1820 she found the first general
result on Fermat's problem: for every prime $p$ with $2p+1$ also prime, there is no solution in which $p$ divides none of
$x,y,z$. Legendre published it in 1825 with credit to her, and it is Exercise~\ref{ex:germain}. Dirichlet and
Legendre settled $n=5$ in 1825, and Lamé settled $n=7$ in 1839 by a computation so heavy that nobody wanted to go to $11$.

Chapter~2 told what happened on 1 March 1847: Lamé's general ``proof'' through the factorisation $x^p+y^p=\prod_k
(x+\zeta^ky)$, Liouville's awkward question about unique factorisation, and Kummer's letter from Breslau showing that it
fails in $\Z[\zeta_{23}]$. Exercise~\ref{ex:kummer23} lets you catch that failure with your own hands. But the story did not
end with Lamé's embarrassment. Kummer did not stop at saying no. Within weeks he had found exactly what Lamé's argument
needs: not that $\Z[\zeta_p]$ has unique factorisation, but that the prime $p$ does not divide its class number. He called
such primes \emph{regular}, proved Fermat's Last Theorem for all of them, and in 1850 he checked the primes below $100$.
Only three failed: $37$, $59$ and $67$. In 1857 the Paris Academy, which had offered a gold medal for a proof of
Fermat's theorem and received nothing it could accept, gave the prize to Kummer, who had not entered the competition.

This chapter retraces Kummer's road with Dedekind's tools. The field $\Q(\zeta_p)$ has a strange and beautiful prime above
$p$, the number $1-\zeta_p$, whose $(p-1)$-st power is $p$ up to a unit. Its ring of integers is as simple as possible,
$\Z[\zeta_n]$. It contains a hidden square root of $\pm p$, the Gauss sum, and with it quadratic reciprocity becomes a
statement about one Frobenius element. Its units are almost real (Kummer's lemma). And for regular primes these pieces fit
together into a proof of the first case of Fermat's Last Theorem. The full theorem had to wait for Andrew Wiles in 1994,
but the road from Kummer to Wiles runs through every chapter of this book: ideals, class groups, $L$-functions, class
field theory, elliptic curves and modular forms.

\begin{reminder}[from Tomes 1, 2 and 4, and Chapters 1--4]
\emph{Cyclotomic polynomials} (Tome~2, Proposition~10.9 and Exercise~10.8). $\Phi_n=\prod_{\gcd(a,n)=1}(X-\zeta_n^a)\in\Z[X]$
is irreducible, $[\Q(\zeta_n):\Q]=\varphi(n)$, and $\operatorname{Gal}(\Q(\zeta_n)/\Q)\cong(\Z/n\Z)^*$ through
$\sigma_a(\zeta_n)=\zeta_n^a$; $X^n-1=\prod_{d\mid n}\Phi_d$. \emph{Chapter~1}: the discriminant of a polynomial,
$\operatorname{disc}(f)=(-1)^{n(n-1)/2}N(f'(\theta))$ (Proposition~\ref{prop:disc}); the index formula (Proposition~\ref{prop:index});
an Eisenstein polynomial at $p$ gives an index prime to $p$ (Proposition~\ref{prop:eisenstein}); the sign of $d_K$ is
$(-1)^{r_2}$ (Proposition~\ref{prop:sign}). \emph{Chapter~2}: unique factorisation of ideals (Theorem~\ref{thm:ufi}),
$\sum e_if_i=n$ (Theorem~\ref{thm:sumef}), the Kummer--Dedekind theorem (Theorem~\ref{thm:kd}), Frobenius elements
(Theorem~\ref{thm:decomp}) and splitting in $\Q(\zeta_m)$ for $p\nmid m$ (Theorem~\ref{thm:cyclosplit}). \emph{Chapter~4}: an
algebraic integer all of whose conjugates have absolute value $1$ is a root of unity (Proposition~\ref{prop:logkernel}).
\end{reminder}

\section{The field of $n$-th roots of unity}

Throughout, $\zeta_n=e^{2i\pi/n}$, and $K_n=\Q(\zeta_n)$. Since $-\zeta_n$ is a primitive $2n$-th root of unity when $n$ is
odd, $K_n=K_{2n}$ for odd $n$; so we may always assume $n\not\equiv2\pmod4$. Apart from $K_1=K_2=\Q$, every $K_n$ is
totally complex of degree $\varphi(n)$, with $r_2=\varphi(n)/2$.

\begin{proposition}\label{prop:muKn}
Let $n\not\equiv2\pmod 4$. The roots of unity in $K_n$ are the $n$-th roots of unity if $n$ is even, and the $2n$-th roots
of unity, $\pm\zeta_n^k$, if $n$ is odd.
\end{proposition}
\begin{proof}
If $\zeta_m\in K_n$, then $K_n$ contains $\zeta_N$ with $N=\operatorname{lcm}(m,n)$ (a product of suitable powers of
$\zeta_m$ and $\zeta_n$, by Bezout), so $\varphi(N)\le\varphi(n)$. Write $N=nt$. Then
\[
\varphi(N)=\varphi(n)\cdot t\prod_{\ell\mid t,\ \ell\nmid n}\Bigl(1-\frac1\ell\Bigr)\ge\varphi(n)\prod_{\ell\mid t,\ \ell\nmid n}(\ell-1),
\]
because $t$ is a multiple of the product of the primes $\ell$ in question. Equality $\varphi(N)=\varphi(n)$ forces
$t=\prod\ell$ with every $\ell-1=1$: so $t=1$, or $t=2$ with $n$ odd. Hence $m\mid n$, or $m\mid2n$ with $n$ odd.
\end{proof}

\begin{figure}[ht]
\centering
\fig{ch05-roots}
\caption{The seventh roots of unity, and the six chords from $1$.}
\label{fig:roots}
\end{figure}

Figure~\ref{fig:roots} shows $K_7$ from the inside. The six chords from $1$ to the other seventh roots of unity have lengths
whose product is $7$: indeed $\prod_{a=1}^{6}(X-\zeta^a)=\Phi_7(X)=1+X+\dots+X^6$, and $X=1$ gives $7$. This tiny identity
controls everything in the next section.

\section{The prime above $p$ and the ring of integers}

\begin{proposition}[the prime $1-\zeta$]\label{prop:lambda}
Let $n=p^k$ be a prime power, $\zeta=\zeta_{p^k}$, $K=\Q(\zeta)$, $\mathcal O=\mathcal O_K$, and $\lambda=1-\zeta$.
\begin{enumerate}[label=(\alph*)]
\item For $p\nmid a$, the number $\dfrac{1-\zeta^a}{1-\zeta}$ is a unit of $\Z[\zeta]$.
\item $\prod_{a\in(\Z/p^k\Z)^*}(1-\zeta^a)=p$, so $p=u\lambda^{\varphi(p^k)}$ with $u$ a unit.
\item $\lambda\mathcal O$ is a prime ideal of norm $p$, and $p\mathcal O=(\lambda\mathcal O)^{\varphi(p^k)}$: the prime $p$ is
\emph{totally ramified} in $K$, with $e=\varphi(p^k)$ and $f=g=1$. Moreover $\lambda\mathcal O\cap\Z=p\Z$ and $\mathcal
O/\lambda\mathcal O=\mathbb F_p$.
\end{enumerate}
\end{proposition}
\begin{proof}
(a) $\frac{1-\zeta^a}{1-\zeta}=1+\zeta+\dots+\zeta^{a-1}\in\Z[\zeta]$. Choose $b\ge1$ with $ab\equiv1\pmod{p^k}$; then
$\zeta=(\zeta^a)^b$, so the inverse $\frac{1-\zeta}{1-\zeta^a}=1+\zeta^a+\dots+\zeta^{a(b-1)}$ is in $\Z[\zeta]$ too.

(b) The roots of $\Phi_p(X^{p^{k-1}})$ are the $x$ with $x^{p^{k-1}}$ a primitive $p$-th root of unity, that is, the primitive
$p^k$-th roots of unity. Both $\Phi_p(X^{p^{k-1}})$ and $\Phi_{p^k}$ are monic of degree $p^{k-1}(p-1)$ with these distinct
roots, so they are equal. Hence $\prod_a(1-\zeta^a)=\Phi_{p^k}(1)=\Phi_p(1)=p$, and by (a) each factor is $\lambda$ times a
unit.

(c) By (b), $p\mathcal O=(\lambda\mathcal O)^{\varphi}$ with $\varphi=\varphi(p^k)=[K:\Q]$. Write $\lambda\mathcal O=\prod
P_i^{a_i}$ (Theorem~\ref{thm:ufi}); then $p\mathcal O=\prod P_i^{\varphi a_i}$, and Theorem~\ref{thm:sumef} gives $\sum
\varphi a_if_i=\varphi$. So there is one $P_1$, with $a_1=f_1=1$: $\lambda\mathcal O=P_1$ is prime with residue field
$\mathbb F_p$. Finally $\lambda\mathcal O\cap\Z$ is a prime ideal of $\Z$ containing $p$, so it is $p\Z$.
\end{proof}

\begin{think}
Total ramification makes congruences modulo $\lambda$ very concrete. Every element of $\mathcal O$ is congruent modulo
$\lambda$ to one of $0,1,\dots,p-1$. Before reading on, find which integer $\zeta^5+3\zeta^2$ is congruent to modulo
$\lambda$ in $\Z[\zeta_7]$. (Since $\zeta\equiv1$, the answer is $4$.)
\end{think}

\begin{theorem}\label{thm:Zzeta}
For every $n$, the ring of integers of $\Q(\zeta_n)$ is $\Z[\zeta_n]$. For a prime power,
\[
d_{\Q(\zeta_{p^k})}=\pm p^{\,p^{k-1}(k(p-1)-1)},\qquad\text{and in particular}\qquad d_{\Q(\zeta_p)}=(-1)^{\frac{p-1}2}p^{\,p-2}
\quad(p\text{ odd}).
\]
\end{theorem}

We prove it first for prime powers, where Eisenstein does all the work, and then glue prime powers together.

\begin{proof}[Proof for $n=p^k$]
Let $\theta=\zeta-1$, so $\Z[\theta]=\Z[\zeta]$, with minimal polynomial $f(X)=\Phi_{p^k}(X+1)$. Modulo $p$, using
$(a+b)^p\equiv a^p+b^p$,
\begin{align*}
f(X)&=\Phi_p\bigl((X+1)^{p^{k-1}}\bigr)\\&\equiv\Phi_p\bigl(X^{p^{k-1}}+1\bigr)\pmod p,\\
\Phi_p(Y+1)&=\frac{(Y+1)^p-1}{Y}\equiv Y^{p-1}\pmod p,
\end{align*}
so $f\equiv X^{\varphi(p^k)}\pmod p$, and its constant term is $f(0)=\Phi_{p^k}(1)=p$. So $f$ is Eisenstein at $p$, and
$p\nmid[\mathcal O:\Z[\zeta]]$ by Proposition~\ref{prop:eisenstein}.

Now compute $\operatorname{disc}(\Phi_{p^k})=\pm N(\Phi_{p^k}'(\zeta))$. Differentiate $X^{p^k}-1=\Phi_{p^k}(X)\bigl(X^{p^{k-1}}
-1\bigr)$ and put $X=\zeta$:
\[
p^k\zeta^{p^k-1}=\Phi'_{p^k}(\zeta)\,(\xi-1),\qquad \xi=\zeta^{p^{k-1}}\text{ a primitive $p$-th root of unity}.
\]
Take norms from $K$ to $\Q$. The norm of $\zeta$ is $\pm1$, and $N(p^k)=p^{k\varphi(p^k)}$. The automorphisms $\sigma_a$ send
$\xi$ to $\xi^a$, and as $a$ runs over $(\Z/p^k\Z)^*$, $a\bmod p$ takes each non-zero value $p^{k-1}$ times; so
$N(\xi-1)=\pm\bigl(\prod_{b=1}^{p-1}(1-\xi^b)\bigr)^{p^{k-1}}=\pm p^{p^{k-1}}$. Hence
$|N(\Phi'_{p^k}(\zeta))|=p^{k\varphi(p^k)-p^{k-1}}=p^{p^{k-1}(k(p-1)-1)}$. By the index formula (Proposition~\ref{prop:index}),
$[\mathcal O:\Z[\zeta]]^2$ divides this power of $p$; being prime to $p$, the index is $1$. So $\mathcal O=\Z[\zeta]$, and $d_K$
is the discriminant just computed. Its sign is $(-1)^{r_2}$ (Proposition~\ref{prop:sign}), with $r_2=\frac{p-1}2$ when
$k=1$.
\end{proof}

For the gluing we need a lemma that is useful far beyond cyclotomic fields.

\begin{lemma}[coprime discriminants]\label{lem:compositum}
Let $K,L$ be number fields of degrees $n,m$ such that $[KL:\Q]=nm$ and $\gcd(d_K,d_L)=1$. If $(\alpha_i)$ and $(\beta_j)$ are
integral bases of $K$ and $L$, then $(\alpha_i\beta_j)$ is an integral basis of $KL$. In particular $\mathcal O_{KL}=\mathcal
O_K\mathcal O_L$.
\end{lemma}
\begin{proof}
The $nm$ products $\alpha_i\beta_j$ span $KL$ over $\Q$ (they span the ring $K[L]$, which is a field since it is finite
dimensional and a domain), so they form a basis. Let $x\in\mathcal O_{KL}$ and write $x=\sum_{i,j}c_{ij}\alpha_i\beta_j$
with $c_{ij}\in\Q$. We show $d_Kc_{ij}\in\Z$.

Let $\theta$ generate $K$. Its minimal polynomial over $L$ has degree $[KL:L]=n=[K:\Q]$, so it is the minimal polynomial over
$\Q$. Hence each of the $n$ embeddings $\sigma_1,\dots,\sigma_n$ of $K$ extends to an embedding of $KL$ that is the identity on
$L$ (send $\theta$ to the corresponding conjugate). Put $x_i=\sum_jc_{ij}\beta_j\in L$; then $\sigma_k(x)=\sum_i\sigma_k(\alpha_i)
x_i$ for $k=1,\dots,n$. Solve this system by Cramer's rule: $x_i=\gamma_i/\delta$, where $\delta=\det(\sigma_k(\alpha_i))$ and
$\gamma_i$ is a determinant whose entries are algebraic integers. So $\delta^2x_i=\delta\gamma_i$ is an algebraic integer.
But $\delta^2=d_K$ (Proposition~\ref{prop:disc}), so $d_Kx_i=\sum_j d_Kc_{ij}\beta_j$ lies in $L\cap\overline{\Z}=\mathcal O_L$, and
since $(\beta_j)$ is an integral basis, $d_Kc_{ij}\in\Z$. By symmetry $d_Lc_{ij}\in\Z$; as $\gcd(d_K,d_L)=1$, $c_{ij}\in\Z$.
\end{proof}

\begin{proof}[Proof of Theorem~\ref{thm:Zzeta} in general]
By induction on the number of prime factors of $n$. Write $n=mp^k$ with $p\nmid m$, $K=\Q(\zeta_m)$ and $L=\Q(\zeta_{p^k})$. Then
$KL=\Q(\zeta_n)$, since $\zeta_m\zeta_{p^k}$ is a primitive $n$-th root of unity and $\zeta_m,\zeta_{p^k}$ are powers of
$\zeta_n$; and $[KL:\Q]=\varphi(n)=\varphi(m)\varphi(p^k)$. By induction $\mathcal O_K=\Z[\zeta_m]$, so $d_K=\pm\operatorname{disc}
(\Phi_m)$, which divides a power of $m$ (proof of Theorem~\ref{thm:cyclosplit}); and $d_L$ is a power of $p$. By the lemma,
$\mathcal O_{KL}=\Z[\zeta_m]\Z[\zeta_{p^k}]=\Z[\zeta_n]$.
\end{proof}

\begin{corollary}\label{cor:cycloram}
Let $n\not\equiv2\pmod4$. A prime $p$ ramifies in $\Q(\zeta_n)$ if and only if $p\mid n$. If $p\nmid n$, it splits into
$\varphi(n)/f$ primes of residue degree $f$, the order of $p$ modulo $n$.
\end{corollary}
\begin{proof}
The case $p\nmid n$ is Theorem~\ref{thm:cyclosplit}. If $p^k\,\|\,n$ with $k\ge1$, then $\varphi(p^k)>1$ (as $p^k\ne2$), and
$L=\Q(\zeta_{p^k})\subseteq\Q(\zeta_n)$. By Proposition~\ref{prop:lambda}, $p\mathcal O_L=(\lambda\mathcal O_L)^{\varphi(p^k)}$, so
$p\mathcal O_{K_n}=(\lambda\mathcal O_{K_n})^{\varphi(p^k)}$, and every prime of $K_n$ above $p$ occurs with exponent at least
$\varphi(p^k)>1$.
\end{proof}

\begin{figure}[ht]
\centering
\fig{ch05-spec}
\caption{The primes of $\Z[\zeta_5]$ above the primes of $\Z$. The fibre above $p$ depends only on $p\bmod 5$: four primes
of degree $1$ when $p\equiv1$, two of degree $2$ when $p\equiv4$, one of degree $4$ when $p\equiv2,3$; and above $5$ the
four sheets of $\operatorname{Spec}\Z[\zeta_5]$ pinch together into the single prime $(1-\zeta)$, counted four times.}
\label{fig:spec5}
\end{figure}

\section{Gauss sums and quadratic reciprocity}

Let $p$ be an odd prime, $\zeta=\zeta_p$, $K=\Q(\zeta)$, $G=\operatorname{Gal}(K/\Q)\cong(\Z/p\Z)^*$, and
$p^*=\bigl(\frac{-1}p\bigr)p=(-1)^{(p-1)/2}p$. Since $G$ is cyclic of even order $p-1$, it has exactly one subgroup of index
$2$, the squares; so $K$ has exactly one quadratic subfield. Which one? The answer is a sum whose sign cost Gauss four years of work.

\begin{proposition}[the Gauss sum]\label{prop:gausssum}
Let $g=\sum_{a=1}^{p-1}\bigl(\frac ap\bigr)\zeta^a$. Then $\sigma_b(g)=\bigl(\frac bp\bigr)g$ for all $b$, and $g^2=p^*$.
Hence the quadratic subfield of $\Q(\zeta_p)$ is $\Q(\sqrt{p^*})$.
\end{proposition}
\begin{proof}
$\sigma_b(g)=\sum_a\bigl(\frac ap\bigr)\zeta^{ab}=\bigl(\frac bp\bigr)\sum_a\bigl(\frac{ab}p\bigr)\zeta^{ab}=\bigl(\frac bp
\bigr)g$, since $ab$ runs over $(\Z/p\Z)^*$ with $a$. Next, $g^2=\sum_{a,b}\bigl(\frac{ab}p\bigr)\zeta^{a+b}$; put $b=ac$
(for fixed $a$, $c$ runs over $(\Z/p\Z)^*$ with $b$), and use $\bigl(\frac{a^2c}p\bigr)=\bigl(\frac cp\bigr)$:
\[
g^2=\sum_{c=1}^{p-1}\Bigl(\frac cp\Bigr)\sum_{a=1}^{p-1}\zeta^{a(1+c)}.
\]
The inner sum is $p-1$ if $c\equiv-1$, and $-1$ otherwise (the sum of all non-trivial $p$-th roots of unity). Since
$\sum_c\bigl(\frac cp\bigr)=0$, we get $g^2=\bigl(\frac{-1}p\bigr)(p-1)+\bigl(\frac{-1}p\bigr)=p^*$. Finally $g\notin\Q$
since $p^*$ is not a rational square, so $\Q(g)=\Q(\sqrt{p^*})$ is the quadratic subfield.
\end{proof}

So $\sqrt5=1+2\zeta_5+2\zeta_5^4$ (check it numerically!), $\sqrt{-7}=\zeta_7+\zeta_7^2+\zeta_7^4-\zeta_7^3-\zeta_7^5-\zeta_7^6$,
and every quadratic field lives inside a cyclotomic field (Exercise~\ref{ex:quadcyclo}). Tome~1, Exercise~6.7 used the Gauss
sum to prove quadratic reciprocity by a clever computation. Now we can see what the computation meant: reciprocity is one
Frobenius element looked at from two sides.

\begin{theorem}[quadratic reciprocity, through Frobenius]\label{thm:qrfrob}
For distinct odd primes $p,q$, $\displaystyle\Bigl(\frac{p^*}q\Bigr)=\Bigl(\frac qp\Bigr)$.
\end{theorem}
\begin{proof}
Let $F=\Q(g)=\Q(\sqrt{p^*})$, and ask whether $q$ splits in $F$. We answer twice.

\emph{From below.} $g$ is an algebraic integer with minimal polynomial $X^2-p^*$, and $[\mathcal O_F:\Z[g]]$ is $1$ or $2$
(Tome~1, Proposition~11.2), hence prime to $q$. By Kummer--Dedekind (Theorem~\ref{thm:kd}), $q$ splits in $F$ if and only if
$X^2-p^*$ has two distinct roots modulo $q$, that is, iff $\bigl(\frac{p^*}q\bigr)=1$.

\emph{From above.} Let $P$ be a prime of $\mathcal O_K$ above $q$ and $\mathfrak q=P\cap\mathcal O_F$. By
Theorem~\ref{thm:cyclosplit}, $\operatorname{Frob}_P=\sigma_q$, and $\sigma_q(x)\equiv x^q\pmod P$ for $x\in\mathcal O_K$.
For $x\in\mathcal O_F$ both sides lie in $\mathcal O_F$ ($\sigma_q$ maps $F$ to itself, $G$ being abelian), so
$\sigma_q(x)\equiv x^q\pmod{\mathfrak q}$. If $\sigma_q$ is the identity on $F$, then $x^q\equiv x$ for every $x$ in the
residue field $\mathcal O_F/\mathfrak q$, which therefore has at most $q$ elements: $f=1$, and since $q$ is unramified in $F$
($q\nmid p^*=d_F$, Theorem~\ref{thm:ramdisc}), $q$ splits. If $\sigma_q$ is not the identity on $F$, then $\sigma_q(g)=-g$,
so $g^q\equiv-g\not\equiv g\pmod{\mathfrak q}$, because $2g\notin\mathfrak q$ (its square $4p^*$ is prime to $q$). So the
residue field is not $\mathbb F_q$, and $q$ is inert. Hence $q$ splits in $F$ iff $\sigma_q(g)=g$, iff $\bigl(\frac qp\bigr)=1$ by
Proposition~\ref{prop:gausssum}.
\end{proof}

\begin{think}
The proof never computed anything modulo $q$ except in $F$. Where exactly did the ``reciprocity'' (the exchange of the roles
of $p$ and $q$) happen? Answer: in the identification $\operatorname{Frob}=\sigma_q$, which says that the Frobenius at $q$
is \emph{the number $q$ itself}, read modulo $p$ in $(\Z/p\Z)^*$. Class field theory (Chapter~10) is the statement that this
miracle happens in every abelian extension.
\end{think}

\section{Units: Kummer's lemma}

Let $p$ be an odd prime, $K=\Q(\zeta_p)$, and $K^+=\Q(\zeta+\zeta^{-1})$. Since $\zeta$ is a root of $X^2-(\zeta+\zeta^{-1})X
+1$, $[K:K^+]\le2$; and $K^+\subseteq\R$ while $K\not\subseteq\R$, so $[K:K^+]=2$ and $K^+=K\cap\R$, the \emph{maximal real
subfield}, of degree $\frac{p-1}2$. It is the fixed field of complex conjugation $\sigma_{-1}$, and since $G$ is abelian,
complex conjugation commutes with every embedding of $K$: $\overline{\sigma(x)}=\sigma(\bar x)$.

Count units with Dirichlet's theorem (Theorem~\ref{thm:dirichlet}). $K$ is totally complex with $r_2=\frac{p-1}2$, so its unit
rank is $\frac{p-3}2$. $K^+$ is totally real of degree $\frac{p-1}2$, so its unit rank is also $\frac{p-3}2$. The real units
already fill the whole rank! Kummer's lemma says much more: they fill almost everything.

\begin{theorem}[Kummer's lemma]\label{thm:kummerlemma}
Every unit $u$ of $\Z[\zeta_p]$ is of the form $u=\zeta^k\varepsilon$ with $k\in\Z$ and $\varepsilon$ a real unit, $\varepsilon\in
\mathcal O_{K^+}^\times$.
\end{theorem}
\begin{proof}
Let $v=u/\bar u$, a unit of $\mathcal O_K$. For every embedding $\sigma$, $|\sigma(v)|=|\sigma(u)|/|\overline{\sigma(u)}|=1$.
By Proposition~\ref{prop:logkernel}, $v$ is a root of unity, so $v=\pm\zeta^j$ by Proposition~\ref{prop:muKn}.

We show the sign is $+$. Every $x\in\mathcal O$ is congruent to an integer modulo $\lambda=1-\zeta$
(Proposition~\ref{prop:lambda}(c)). Also $\bar\lambda=1-\zeta^{-1}=-\zeta^{-1}\lambda$, so $\bar\lambda\mathcal O=\lambda\mathcal
O$ and conjugation preserves congruences modulo $\lambda$: if $x\equiv b\pmod\lambda$ with $b\in\Z$, then $\bar x\equiv b
\pmod\lambda$ too. Choose $\eta=\zeta^{j(p+1)/2}$, so that $\eta^2=\zeta^j$, and let $w=u\eta^{-1}$; then $\bar w=\bar u\eta$. If
$u=-\zeta^j\bar u$, then $w=-\eta\bar u=-\bar w$. Writing $w\equiv b\pmod\lambda$ with $b\in\Z$, we get $b\equiv-b$, so
$\lambda\mid2b$, so $p\mid 2b$ (Proposition~\ref{prop:lambda}(c)), so $p\mid b$ and $\lambda\mid w$: impossible for a unit.

Hence $u=\zeta^j\bar u$, that is $w=\bar w$: $\varepsilon=w$ is a real unit and $u=\eta\varepsilon$ with $\eta$ a power of
$\zeta$.
\end{proof}

The real units contain explicit specimens, the \emph{cyclotomic units}. By Proposition~\ref{prop:lambda}(a),
$\frac{1-\zeta^a}{1-\zeta}$ is a unit, and multiplying by the right power of $\zeta$ makes it real: up to sign it equals
\[
\xi_a=\frac{\sin(\pi a/p)}{\sin(\pi/p)},\qquad 1<a<\frac p2.
\]
For $p=5$, $\xi_2=2\cos\frac\pi5$ is the golden ratio (Exercise~\ref{ex:golden}); for $p=7$, $\xi_2$ and $\xi_3$ are, up to sign,
the units $\theta$ and $1+\theta$ of Chapter~4, Figure~\ref{fig:loglattice}.

\begin{remark}
Kummer proved that the $\xi_a$ generate a subgroup of $\mathcal O_{K^+}^\times$ of finite index, and that this index is the
class number $h^+$ of $K^+$. The independence of the $\xi_a$ amounts to the non-vanishing of $L(1,\chi)$ for the even
Dirichlet characters modulo $p$, which Chapter~8 proves; the equality with $h^+$ is in Washington's book, Theorem~8.2.
\end{remark}

\section{Regular primes and Fermat's Last Theorem}

\begin{definition}
An odd prime $p$ is \emph{regular} if $p$ does not divide the class number $h_p$ of $\Q(\zeta_p)$, and \emph{irregular}
otherwise.
\end{definition}

For $p\le19$, $\Z[\zeta_p]$ is a principal ideal domain ($h_p=1$), so these primes are regular. For $p=23$, $h_p=3$: the ring
is not a UFD, which is what Kummer had found, but $23$ is still regular. How can one decide regularity without computing
class groups of fields of degree $p-1$? Kummer found a stunning criterion.

\begin{remark}[Kummer's criterion]
\emph{$p$ is irregular if and only if $p$ divides the numerator of one of the Bernoulli numbers $B_2,B_4,\dots,B_{p-3}$}
(Tome~4, Proposition~1.14). Chapter~8 proves, through the analytic class number formula, the half we need to recognise an
irregular prime: if $p$ divides such a numerator, then $p\mid h_p$. The converse is Theorem~5.34 in Washington's book. With it,
the irregular primes below $160$ are
\begin{gather*}
37\ (B_{32}),\quad 59\ (B_{44}),\quad 67\ (B_{58}),\\ 101\ (B_{68}),\quad 103\ (B_{24}),\quad 131\ (B_{22}),\\
149\ (B_{130}),\quad 157\ (B_{62},B_{110}).
\end{gather*}
Jensen proved in 1915 that there are infinitely many irregular primes. Nobody knows whether there are infinitely many regular
ones, although about $e^{-1/2}\approx61\%$ of all primes should be regular (a numerical calculation tests this).
\end{remark}

Here is why regularity is exactly what Lamé needed.

\begin{lemma}\label{lem:regular}
Let $p$ be regular and $I$ an ideal of $\Z[\zeta_p]$. If $I^p$ is principal, then $I$ is principal.
\end{lemma}
\begin{proof}
The class $[I]$ satisfies $[I]^p=1$ and $[I]^{h_p}=1$. Since $\gcd(p,h_p)=1$, write $1=ap+bh_p$; then
$[I]=[I]^{ap}[I]^{bh_p}=1$.
\end{proof}

\begin{think}
Lamé's argument says: the factors $x+\zeta^ky$ are pairwise coprime and their product is a $p$-th power, so each is a $p$-th
power. Which two steps of this sentence fail in a general $\Z[\zeta_p]$, and which of them does unique factorisation of
\emph{ideals} rescue? (Coprimality makes sense for ideals, and each ideal $(x+\zeta^ky)$ is the $p$-th power of an ideal.
What is lost is that this ideal is principal, which Lemma~\ref{lem:regular} restores; and what Lamé forgot entirely is the
unit in front, which Kummer's lemma tames.)
\end{think}

Two small lemmas about congruences modulo $p$ in $\Z[\zeta]$, where ``$\alpha\equiv\beta\pmod p$'' means $\alpha-\beta\in p\Z
[\zeta]$.

\begin{lemma}\label{lem:pthpower}
For every $\alpha\in\Z[\zeta_p]$ there is $a\in\Z$ with $\alpha^p\equiv a\pmod p$.
\end{lemma}
\begin{proof}
Write $\alpha=\sum_ia_i\zeta^i$. In the ring $\Z[\zeta]/p\Z[\zeta]$ of characteristic $p$, the map $x\mapsto x^p$ is additive, so
$\alpha^p\equiv\sum_ia_i^p\zeta^{ip}=\sum_ia_i^p\pmod p$.
\end{proof}

\begin{lemma}\label{lem:zerocoef}
Let $\alpha=\sum_{i=0}^{p-1}a_i\zeta^i$ with $a_i\in\Z$, at least one $a_i=0$. If $n\in\Z$ and $\alpha\equiv0\pmod n$, then
$n$ divides every $a_i$.
\end{lemma}
\begin{proof}
$1,\zeta,\dots,\zeta^{p-2}$ is a $\Z$-basis of $\Z[\zeta]$ (Theorem~\ref{thm:Zzeta}), and since $\zeta^{p-1}=-1-\zeta-\dots-
\zeta^{p-2}$, any $p-1$ of the $p$ numbers $1,\zeta,\dots,\zeta^{p-1}$ form a $\Z$-basis. If $a_k=0$, then $\alpha$ is written
in the basis that omits $\zeta^k$, with coefficients $a_i$; and $\alpha/n\in\Z[\zeta]$ means that all these coefficients are
divisible by $n$.
\end{proof}

We can now prove Kummer's theorem in the ``first case'', where $p$ divides none of $x,y,z$. This is the case Lamé's
argument was aiming at, and it is where the beauty is.

\begin{theorem}[Kummer, 1847]\label{thm:fltfirst}
Let $p\ge5$ be a regular prime. There are no integers $x,y,z$ with $x^p+y^p=z^p$ and $p\nmid xyz$.
\end{theorem}
\begin{proof}
Suppose there are. Dividing by a common factor, we may assume $x,y,z$ pairwise coprime (a prime dividing two of them divides
the third).

\emph{Step 0: we may assume $x\not\equiv y\pmod p$.} If $x\equiv y\equiv-z\pmod p$, then $z^p=x^p+y^p$ gives $z\equiv x+y
\equiv-2z$, so $p\mid3z$, impossible as $p\ge5$ and $p\nmid z$. So either $x\not\equiv y$, or $x\not\equiv-z$; in the second
case rewrite the equation as $x^p+(-z)^p=(-y)^p$ ($p$ is odd) and rename.

\emph{Step 1: the ideals $(x+\zeta^iy)$, $0\le i<p$, are pairwise coprime.} Let $P$ be a prime containing $x+\zeta^iy$ and
$x+\zeta^jy$ with $i\ne j$. Then $P$ contains their difference $(\zeta^i-\zeta^j)y$ and $\zeta^j(x+\zeta^iy)-\zeta^i(x+
\zeta^jy)=(\zeta^j-\zeta^i)x$. Now $\zeta^i-\zeta^j$ is $\lambda$ times a unit (Proposition~\ref{prop:lambda}(a)), so $P$
contains $\lambda x$ and $\lambda y$, hence $\lambda$ (as $ux+vy=1$ for some integers $u,v$). So $P=(\lambda)$. But then
$x+y\equiv x+\zeta^iy\equiv0\pmod\lambda$, so $p\mid x+y$, and $z^p=x^p+y^p\equiv x+y\equiv0\pmod p$: $p\mid z$, a
contradiction.

\emph{Step 2: $x+\zeta y=u\alpha^p$ with $u$ a unit.} The ideal equation $\prod_{i=0}^{p-1}(x+\zeta^iy)=(z)^p$ and unique
factorisation into primes (Theorem~\ref{thm:ufi}), together with Step~1, show that each $(x+\zeta^iy)$ is the $p$-th power of an
ideal: every prime appears in exactly one factor, with an exponent divisible by $p$. So $(x+\zeta y)=I^p$, $I$ is principal
by Lemma~\ref{lem:regular}, say $I=(\alpha)$, and $x+\zeta y=u\alpha^p$.

\emph{Step 3: a congruence.} By Kummer's lemma, $u=\zeta^r\varepsilon$ with $\varepsilon$ real, and by
Lemma~\ref{lem:pthpower}, $\alpha^p\equiv a\pmod p$ with $a\in\Z$. So $x+\zeta y\equiv\zeta^r\varepsilon a\pmod p$. Complex
conjugation maps $p\Z[\zeta]$ to itself and fixes $\varepsilon$ and $a$, so $x+\zeta^{-1}y\equiv\zeta^{-r}\varepsilon a\pmod
p$. Eliminating $\varepsilon a$,
\[
\zeta^{-r}(x+\zeta y)\equiv\zeta^r(x+\zeta^{-1}y),\qquad\text{that is}\qquad
x+\zeta y-\zeta^{2r}x-\zeta^{2r-1}y\equiv0\pmod p.\tag{$*$}
\]

\emph{Step 4: the exponents collide.} Reduce the exponents $0,1,2r,2r-1$ modulo $p$. If they are four distinct residues,
then, as $p\ge5$, some power of $\zeta$ among $1,\dots,\zeta^{p-1}$ does not occur in $(*)$, and Lemma~\ref{lem:zerocoef} gives
$p\mid x$: contradiction. Since $0\ne1$ and $2r\ne2r-1$, the possible coincidences are:
\begin{itemize}
\item $2r\equiv0$. Then $(*)$ reads $y\zeta-y\zeta^{p-1}\equiv0$, and Lemma~\ref{lem:zerocoef} gives $p\mid y$: contradiction.
\item $2r\equiv1$ (equivalently $2r-1\equiv0$). Then $(*)$ reads $(x-y)-(x-y)\zeta\equiv0$, so $p\mid x-y$, which Step~0 excluded.
\item $2r-1\equiv1$. Then $(*)$ reads $x-x\zeta^2\equiv0$, so $p\mid x$: contradiction.
\end{itemize}
Every case is impossible, and the theorem is proved.
\end{proof}

For $p=3$ the first case is a one-line congruence (Exercise~\ref{ex:cubes9}), and Tome~1, Theorem~10.8 settled $n=3$
completely. Since every $n\ge3$ is divisible by $4$ or by an odd prime, and Fermat did $n=4$, Fermat's Last Theorem for all
$n$ reduces to odd prime exponents.

\begin{remark}[the second case, and the end of the story]
Kummer also proved the second case ($p\mid xyz$) for regular primes, by a more delicate descent in $\Z[\zeta_p]$ that
generalises Tome~1, Theorem~10.8; see Washington, Theorem~9.3. For irregular primes he and his successors developed finer
criteria, and by 1993 computers had verified the theorem for all exponents up to four million. Then Andrew Wiles, with a
final step completed with Richard Taylor, proved in 1994 that every semistable elliptic curve over $\Q$ is modular, which by
the work of Frey, Serre and Ribet implies Fermat's Last Theorem for every exponent. The elliptic curves are in
Chapters~12--13 and the modular forms in Chapter~14; the bridge between them is still, thirty years later, one of the most
active subjects in mathematics.
\end{remark}

\section{Exercises}

\begin{exercise}[$\star$]\label{ex:cubes9}
\emph{The first case for $p=3$.} Show that every cube is congruent to $0$ or $\pm1$ modulo $9$, and deduce that $x^3+y^3=z^3$
has no solution with $3\nmid xyz$. Where does the proof of Theorem~\ref{thm:fltfirst} use $p\ge5$?
\end{exercise}

\begin{exercise}[$\star\star$]\label{ex:golden}
\emph{The golden ratio is a cyclotomic unit.} Let $\zeta=\zeta_5$ and $\phi=\frac{1+\sqrt5}2$.
\begin{enumerate}[label=(\alph*)]
\item Show that $\theta=\zeta+\zeta^{-1}$ satisfies $\theta^2+\theta-1=0$, so $\theta=\phi^{-1}$ and $K^+=\Q(\sqrt5)$.
\item Show that $\phi=-\zeta^2(1+\zeta)$, and that $1+\zeta$ is a cyclotomic unit.
\item Show that the units of $\Z[\zeta_5]$ are exactly the numbers $\pm\zeta^k\phi^m$, $k,m\in\Z$.
\end{enumerate}
\end{exercise}

\begin{exercise}[$\star\star$]\label{ex:realint}
\emph{The integers of the real subfield.} Let $p$ be an odd prime, $K=\Q(\zeta_p)$ and $\theta=\zeta+\zeta^{-1}$. Show that
$\mathcal O_{K^+}=\Z[\theta]$.
\end{exercise}

\begin{exercise}[$\star\star$]\label{ex:compdisc}
\emph{Discriminants of composita.} Under the hypotheses of Lemma~\ref{lem:compositum}, show that $d_{KL}=d_K^{\,m}d_L^{\,n}$.
Deduce the discriminant of $\Q(\zeta_{12})=\Q(i,\sqrt{-3})$, and check Theorem~\ref{thm:Zzeta} for $\Q(\zeta_8)$ against
$\operatorname{disc}(X^4+1)$.
\end{exercise}

\begin{exercise}[$\star\star$]\label{ex:quadcyclo}
\emph{A first glimpse of Kronecker--Weber.} Show that every quadratic field $\Q(\sqrt d)$, $d$ a squarefree integer, is
contained in $\Q(\zeta_{4|d|})$.
\end{exercise}

\begin{exercise}[$\star\star\star$]\label{ex:germain}
\emph{Sophie Germain's theorem.} Let $p$ be an odd prime such that $q=2p+1$ is prime (for instance $p=3,5,11,23,29,41$).
Show that $x^p+y^p=z^p$ has no solution in integers with $p\nmid xyz$. (Work with $x^p+y^p+z^p=0$, pairwise coprime.)
\end{exercise}

\begin{exercise}[$\star\star\star$]\label{ex:kummer23}
\emph{Kummer's counterexample.} Let $K=\Q(\zeta_{23})$ and $F=\Q(\sqrt{-23})$.
\begin{enumerate}[label=(\alph*)]
\item Show that $F\subseteq K$ and that $47$ splits completely in $K$.
\item For $\alpha\in\mathcal O_K$, put $N_{K/F}(\alpha)=\prod_{\sigma\in H}\sigma(\alpha)$, where $H=\operatorname{Gal}(K/F)$.
Show that $N_{K/F}(\alpha)\in\mathcal O_F$ and $N_{K/\Q}(\alpha)=N_{F/\Q}(N_{K/F}(\alpha))$.
\item Show that no element of $\mathcal O_F$ has norm $47$, and deduce that no prime of $K$ above $47$ is principal. So
$\Z[\zeta_{23}]$ is not a unique factorisation domain.
\end{enumerate}
\end{exercise}

\begin{beauty}
Count what one root of unity bought us. The single identity $\prod(1-\zeta^a)=p$ made $p$ totally ramified, gave the ring of
integers through Eisenstein, and produced the cyclotomic units. The Gauss sum hid $\sqrt{\pm p}$ inside $\Q(\zeta_p)$, and
the Frobenius $\sigma_q$ turned quadratic reciprocity into the statement that a symmetry of $\Q(\zeta_p)$ restricts to a
symmetry of $\Q(\sqrt{p^*})$. And Kummer's attack on Fermat, born from Lamé's mistake, created ideal theory, class
numbers, and the link between class numbers and Bernoulli numbers, the first appearance of the mysterious bond between
algebra and the values of zeta functions that Chapter~8 starts to explain and that modern Iwasawa theory still explores.
Fermat's margin was too small; the theory it provoked fills libraries.
\end{beauty}

\section{Hints}

Try each exercise for a while before reading its hint, and read the hint before the solution.

\begin{hint}{ex:cubes9}
Check $x=0$, $\pm1$, $\pm2$, $\pm3$, $\pm4$.
\end{hint}
\begin{hint}{ex:golden}
Divide $1+\zeta+\zeta^2+\zeta^3+\zeta^4=0$ by $\zeta^2$. For (c), use Kummer's lemma and the units of $\Q(\sqrt5)$ (Tome~1,
Exercise~11.2).
\end{hint}
\begin{hint}{ex:realint}
$\mathcal O_{K^+}=\Z[\zeta]\cap\R$. Write a real element in the basis $\zeta,\zeta^2,\dots,\zeta^{p-1}$ and compare with its
conjugate.
\end{hint}
\begin{hint}{ex:compdisc}
The embeddings of $KL$ are the pairs (embedding of $K$, embedding of $L$); the matrix $(\sigma\tau(\alpha_i\beta_j))$ is a
Kronecker product.
\end{hint}
\begin{hint}{ex:quadcyclo}
$\sqrt{p^*}\in\Q(\zeta_p)$, $i\in\Q(\zeta_4)$, and $\sqrt2=\zeta_8+\zeta_8^{-1}$.
\end{hint}
\begin{hint}{ex:germain}
The $p$-th powers modulo $q$ are $0$ and $\pm1$. Show $q\mid xyz$, say $q\mid x$, and factor $-x^p=(y+z)(y^{p-1}-\dots+z^{p-1})$
into coprime factors, and similarly for $y$ and $z$.
\end{hint}
\begin{hint}{ex:kummer23}
$23\equiv3\pmod4$, and $47=2\cdot23+1$. Elements of $\mathcal O_F$ are $\frac{a+b\sqrt{-23}}2$ with $a\equiv b\pmod2$.
\end{hint}

\section{Solutions}

\begin{solution}{ex:cubes9}
\idea Modulo $9$, cubes take very few values, because $(3k\pm1)^3\equiv\pm1$.
\writeup The cubes of $0$, $\pm1$, $\pm2$, $\pm3$, $\pm4$ are $0$, $\pm1$, $\pm8$, $\pm27$, $\pm64$, that is $0$, $\pm1$, $\mp1$, $0$,
$\pm1$ modulo $9$. If $3\nmid
xyz$, then $x^3+y^3-z^3\equiv\pm1\pm1\pm1\pmod9$, which is one of $\pm1,\pm3$, never $0$. The proof of
Theorem~\ref{thm:fltfirst} uses $p\ge5$ twice: in Step~0, to rule out $p\mid3z$; and in Step~4, to be sure that four exponents
leave a power of $\zeta$ unused among $p\ge5$ of them. For $p=3$ both steps fail, which is why a separate argument is needed.
\end{solution}

\begin{solution}{ex:golden}
\idea Symmetric expressions in $\zeta$ and $\zeta^{-1}$ are polynomials in $\theta$, and the relation $\Phi_5(\zeta)=0$ becomes a
quadratic equation for $\theta$.
\writeup (a) Dividing $1+\zeta+\zeta^2+\zeta^3+\zeta^4=0$ by $\zeta^2$ gives $(\zeta^2+\zeta^{-2})+(\zeta+\zeta^{-1})+1=0$, and
$\zeta^2+\zeta^{-2}=\theta^2-2$, so $\theta^2+\theta-1=0$. As $\theta=2\cos\frac{2\pi}5>0$, $\theta=\frac{-1+\sqrt5}2=\phi-1=
\phi^{-1}$. So $K^+=\Q(\theta)=\Q(\sqrt5)$.

(b) $-\zeta^2(1+\zeta)=-(\zeta^2+\zeta^3)=-(\zeta^2+\zeta^{-2})=2-\theta^2=2-(1-\theta)=1+\theta=1+\phi^{-1}=\phi$. And
$1+\zeta=\frac{1-\zeta^2}{1-\zeta}$ is a unit by Proposition~\ref{prop:lambda}(a).

(c) By Kummer's lemma, a unit is $\zeta^k\varepsilon$ with $\varepsilon$ a unit of $\mathcal O_{K^+}=\mathcal O_{\Q(\sqrt5)}$,
whose units are $\pm\phi^m$ (Tome~1, Exercise~11.2 and Theorem~11.4). Conversely every $\pm\zeta^k\phi^m$ is a unit. So
$\mathcal O_K^\times=\mu_{10}\times\phi^{\Z}$, a free part of rank $1=\frac{5-3}2$, as Dirichlet predicts.
\end{solution}

\begin{solution}{ex:realint}
\idea A real element equals its conjugate, which pairs up the coefficients of $\zeta^i$ and $\zeta^{-i}$; and
$\zeta^i+\zeta^{-i}$ is a polynomial in $\theta$, as $2\cos(ix)$ is a polynomial in $2\cos x$.
\writeup $K^+=K\cap\R$ (shown in the text), so $\mathcal O_{K^+}=\mathcal O_K\cap\R=\Z[\zeta]\cap\R$, and clearly $\Z[\theta]
\subseteq\mathcal O_{K^+}$. The numbers $\zeta,\dots,\zeta^{p-1}$ form a $\Z$-basis of $\Z[\zeta]$ (they are $\zeta$ times the basis
$1,\dots,\zeta^{p-2}$, and $\zeta$ is a unit). Let $x=\sum_{i=1}^{p-1}a_i\zeta^i$ be real. Then $x=\bar x=\sum_ia_i\zeta^{p-i}$, and
comparing coefficients, $a_i=a_{p-i}$. Hence $x=\sum_{i=1}^{(p-1)/2}a_i(\zeta^i+\zeta^{-i})$. Define $P_0=2$, $P_1=X$,
$P_{i+1}=XP_i-P_{i-1}$; by induction $\zeta^i+\zeta^{-i}=P_i(\theta)$, since $(\zeta+\zeta^{-1})(\zeta^i+\zeta^{-i})=(\zeta^{i+1}+
\zeta^{-i-1})+(\zeta^{i-1}+\zeta^{1-i})$. So $x\in\Z[\theta]$.
\end{solution}

\begin{solution}{ex:compdisc}
\idea The discriminant is the square of the determinant of the matrix of embeddings, and here that matrix splits as a tensor
product.
\writeup Let $\sigma_1,\dots,\sigma_n$ and $\tau_1,\dots,\tau_m$ be the embeddings of $K$ and $L$. An embedding $\rho$ of $KL$ is
determined by its restrictions to $K$ and $L$, which generate $KL$; there are $nm$ embeddings of $KL$ and $nm$ pairs, so every
pair $(\sigma_k,\tau_l)$ comes from exactly one $\rho_{kl}$. By the lemma, $(\alpha_i\beta_j)$ is an integral basis of $KL$, and
$\rho_{kl}(\alpha_i\beta_j)=\sigma_k(\alpha_i)\tau_l(\beta_j)$: the matrix is the Kronecker product $A\otimes B$ of
$A=(\sigma_k(\alpha_i))$ and $B=(\tau_l(\beta_j))$. Since $A\otimes B=(A\otimes I_m)(I_n\otimes B)$, its determinant is
$\det(A)^m\det(B)^n$. Squaring, $d_{KL}=d_K^{\,m}d_L^{\,n}$.

For $\Q(\zeta_{12})$: $K=\Q(\zeta_3)$ with $d_K=-3$ and $L=\Q(i)$ with $d_L=-4$ are coprime, and $[KL:\Q]=4=\varphi(12)$; so
$d=(-3)^2(-4)^2=144$. For $\Q(\zeta_8)$, Theorem~\ref{thm:Zzeta} with $p=2$, $k=3$ gives $\pm2^{4\cdot(3-1)}=\pm2^8$, with sign
$(-1)^{r_2}=+1$. Directly, $\operatorname{disc}(X^4+1)=N(4\zeta^3)=4^4N(\zeta)^3=256$ (Proposition~\ref{prop:disc}(c), with
$(-1)^{4\cdot3/2}=1$).
\end{solution}

\begin{solution}{ex:quadcyclo}
\idea Build $\sqrt d$ as a product of square roots that we already know to be cyclotomic: $\sqrt{p^*}$ from the Gauss sum,
$i$, and $\sqrt2$.
\writeup Write $d=\pm2^{\epsilon}p_1\cdots p_r$ with distinct odd primes $p_i$ and $\epsilon\in\{0,1\}$. Each $\sqrt{p_i^*}$ lies
in $\Q(\zeta_{p_i})$ (Proposition~\ref{prop:gausssum}), $i=\zeta_4$, and $\sqrt2=\zeta_8+\zeta_8^{-1}=2\cos\frac\pi4$. Since
$p_i=\pm p_i^*$, the number $\sqrt d$ is a product of the $\sqrt{p_i^*}$, possibly $i$, and possibly $\sqrt2$. All of these lie
in $\Q(\zeta_N)$ with $N=\operatorname{lcm}(4,|d'|)$ if $d$ is odd ($d'=d$), and $N=\operatorname{lcm}(8,|d|/2)$ if $d$ is
even; in both cases $N$ divides $4|d|$, and $\Q(\zeta_N)\subseteq\Q(\zeta_{4|d|})$. The Kronecker--Weber theorem (Chapter~10)
says that \emph{every} abelian extension of $\Q$ lies in some $\Q(\zeta_N)$.
\end{solution}

\begin{solution}{ex:germain}
\idea Modulo $q=2p+1$, a $p$-th power is a square root of $1$ or $0$, so the equation forces $q$ to divide one variable. Then
the three factorisations $-x^p=(y+z)Q(y,z)$, and its two cousins, turn the variables into $p$-th powers, and modulo $q$ the
numbers become too few to fit.
\writeup Write the equation as $x^p+y^p+z^p=0$ (replace $z$ by $-z$), with $x,y,z$ pairwise coprime and $p\nmid xyz$.

\emph{$p$-th powers modulo $q$.} If $q\nmid t$, then $(t^p)^2=t^{q-1}\equiv1$, so $t^p\equiv\pm1\pmod q$. If $q\nmid xyz$, the
sum of three numbers $\pm1$ would be $\equiv0\pmod q$; but such a sum is odd and of absolute value at most $3<q$. So $q\mid
xyz$; by symmetry, $q\mid x$.

\emph{Factorisations.} Let $Q(y,z)=\sum_{i=0}^{p-1}y^{p-1-i}(-z)^i$, so that $-x^p=y^p+z^p=(y+z)Q(y,z)$. Modulo $y+z$, $z\equiv-y$
and $Q\equiv py^{p-1}$; so a common prime factor of $y+z$ and $Q$ divides $py^{p-1}$. It cannot divide $y$ (it would divide $z$
too), and it cannot be $p$ (it would divide $x$). So $y+z$ and $Q$ are coprime, their product is the $p$-th power $(-x)^p$, and
since $p$ is odd, $y+z=a^p$ and $Q(y,z)=\alpha^p$ for integers $a,\alpha$. In the same way $z+x=b^p$ and $x+y=c^p$.

\emph{Contradiction.} $2x=(x+y)+(z+x)-(y+z)=c^p+b^p-a^p\equiv0\pmod q$, and as before one of $a,b,c$ is divisible by $q$. If $q
\mid b$, then $q\mid z+x$, so $q\mid z$, contradicting $\gcd(x,z)=1$; likewise $q\nmid c$. So $q\mid a$, that is $z\equiv-y
\pmod q$, and $\alpha^p=Q(y,z)\equiv py^{p-1}\pmod q$. But $y\equiv x+y=c^p\pmod q$ with $c^p\equiv\pm1$, so $y^{p-1}\equiv1$ ($p-1$
is even), and $\alpha^p\equiv p\pmod q$. The left side is $0$ or $\pm1$ modulo $q$, while $p,p\pm1$ are not divisible by
$q=2p+1$. This contradiction proves the theorem.
\end{solution}

\begin{solution}{ex:kummer23}
\idea A principal prime $P=(\alpha)$ of norm $47$ would push down, through the relative norm, to an element of norm $47$ in the
quadratic field $\Q(\sqrt{-23})$; and the norm form $\frac{a^2+23b^2}4$ never takes the value $47$.
\writeup (a) $23\equiv3\pmod4$, so $23^*=-23$ and $F=\Q(\sqrt{-23})\subseteq K$ by Proposition~\ref{prop:gausssum}. Since
$47\equiv1\pmod{23}$, the order of $47$ modulo $23$ is $1$, and $47$ splits into $22$ primes of norm $47$
(Corollary~\ref{cor:cycloram}).

(b) $\beta=N_{K/F}(\alpha)$ is a product of algebraic integers, so it is one; and for $\tau\in H$, $\tau(\beta)=\prod_\sigma
\tau\sigma(\alpha)=\beta$, so $\beta\in K^H=F$. Hence $\beta\in\mathcal O_F$. Let $\rho$ be the non-trivial element of
$\operatorname{Gal}(K/\Q)/H$, lifted to $\tilde\rho$. Then $G=H\cup\tilde\rho H$, and
$N_{K/\Q}(\alpha)=\prod_{\sigma\in H}\sigma(\alpha)\prod_{\sigma\in H}\tilde\rho\sigma(\alpha)=\beta\,\tilde\rho(\beta)=N_{F/\Q}(\beta)$.

(c) $\mathcal O_F=\Z\bigl[\frac{1+\sqrt{-23}}2\bigr]$ consists of the $\frac{a+b\sqrt{-23}}2$ with $a\equiv b\pmod2$, of norm
$\frac{a^2+23b^2}4$. Norm $47$ means $a^2+23b^2=188$; $b=0,\pm1,\pm2$ give $a^2=188,165,96$, none a square, and $|b|\ge3$ gives
$23b^2>188$. Now let $P$ be a prime of $K$ above $47$, of norm $47$ by (a). If $P=(\alpha)$, then $N_{K/\Q}(\alpha)=\pm47$, and
in fact $+47$ since $K$ is totally complex (the conjugates come in pairs $\sigma(\alpha),\overline{\sigma(\alpha)}$). By (b),
$\beta=N_{K/F}(\alpha)\in\mathcal O_F$ has norm $47$, which is impossible. So $P$ is not principal, the class number of
$\Q(\zeta_{23})$ is not $1$, and $\Z[\zeta_{23}]$ is not a UFD. (A Dedekind domain that is a UFD is a PID: a prime $P\ne0$ contains an
irreducible factor $\pi$ of any of its non-zero elements; in a UFD $\pi$ is prime, so $(\pi)\subseteq P$ is an inclusion of
non-zero primes, which are maximal, hence an equality.) In fact $h(\Q(\zeta_{23}))=3$.
\end{solution}

\section*{Sources and further reading}

\begin{itemize}
\item E.~E.~Kummer, ``Beweis des Fermat'schen Satzes der Unm\"oglichkeit von $x^\lambda+y^\lambda=z^\lambda$ f\"ur eine unendliche
Anzahl Primzahlen $\lambda$'', \emph{Monatsber.\ Akad.\ Wiss.\ Berlin} (1847), 132--139; and \emph{J.\ reine angew.\ Math.}\ 40
(1850), 130--138. Regular primes and the first proofs.
\item L.~C.~Washington, \emph{Introduction to Cyclotomic Fields}, 2nd ed., Springer GTM~83, 1997. Chapter~1 (the first case, the
route followed here), Chapters~2, 5 and~8 (integers, Kummer's criterion, cyclotomic units), Theorem~9.3 (the second case).
\item H.~M.~Edwards, \emph{Fermat's Last Theorem: A Genetic Introduction to Algebraic Number Theory}, Springer GTM~50, 1977. The
history of Lamé, Liouville and Kummer told from the original sources.
\item P.~Ribenboim, \emph{13 Lectures on Fermat's Last Theorem}, Springer, 1979; and \emph{Fermat's Last Theorem for Amateurs},
Springer, 1999. Sophie Germain's theorem and many elementary results.
\item D.~A.~Marcus, \emph{Number Fields}, 2nd ed., Springer, 2018, Chapters~2 and~3. Cyclotomic integers and the compositum lemma.
\item K.~Ireland and M.~Rosen, \emph{A Classical Introduction to Modern Number Theory}, 2nd ed., Springer GTM~84, 1990,
Chapters~6 and~14. Gauss sums and the Stickelberger relation.
\item A.~Wiles, ``Modular elliptic curves and Fermat's Last Theorem'', \emph{Ann.\ of Math.}\ 141 (1995), 443--551; R.~Taylor and
A.~Wiles, ``Ring-theoretic properties of certain Hecke algebras'', same volume, 553--572.
\item \emph{History.} Fermat's marginal note is translated from the Latin of the 1670 edition of Diophantus (Edwards, above,
\S1.1). Sophie Germain's correspondence with Gauss and her theorem: R.~Laubenbacher and D.~Pengelley, ``Voici ce que j'ai
trouv\'e: Sophie Germain's grand plan to prove Fermat's Last Theorem'', \emph{Historia Math.}\ 37 (2010), 641--692.
The 1857 prize to Kummer and the three irregular primes below $100$: Edwards, above, Chapters~4--6.
\end{itemize}
